\newcolumntype{P}[1]{>{\centering\arraybackslash}p{#1}}
\setlength\tabcolsep{1pt}
\renewcommand{\arraystretch}{1.6}
\newcommand{\bangle}[2]{\genfrac{\langle}{\rangle}{0pt}{}{#1}{#2}}
\newcommand{\stirlingone}[2]{\genfrac{\lbrack}{\rbrack}{0pt}{}{#1}{#2}}
\newcommand{\stirlingtwo}[2]{\genfrac{\lbrace}{\rbrace}{0pt}{}{#1}{#2}}
\providecommand{\myfloor}[1]{ \left\lfloor #1 \right\rfloor }
\begin{footnotesize}

\subsubsection{Mobius Inversion}
	\[F(n)=\sum_{d|n}f(d)\Rightarrow f(n)=\sum_{d|n}\mu(d)F\left(\frac{n}{d}\right)\]
	\[F(n)=\sum_{\substack{n|d\\d\le N}}f(d)
		\Rightarrow f(n)=\sum_{\substack{n|d\\d\le N}}\mu\left(\frac{d}{n}\right)F(d)\]
	\[\relax[x=1]=\sum_{d|x}\mu(d),\qquad x=\sum_{d|x}\varphi(d)\quad(x\ge1)\]
\subsubsection{杜教筛}
	\[
S_f(n)=\sum_{i=1}^n f(i),\qquad
S_{\varphi}(n)=\frac {n(n+1)}2-\sum_{d=2}^n
S_{\varphi}\left(\left\lfloor\frac nd\right\rfloor\right)
\]
\[
S_\mu(n)=1-\sum_{d=2}^n
S_\mu\left(\left\lfloor\frac nd\right\rfloor\right)
\]
\[
q=\left\lfloor\frac nl\right\rfloor,\quad
r=\left\lfloor\frac nq\right\rfloor
\Rightarrow
\sum_{d=l}^{r}S_f\left(\left\lfloor\frac nd\right\rfloor\right)
=(r-l+1)S_f(q).
\]

\subsubsection{降幂公式}
	\[\gcd(a,m)=1\Rightarrow
		a^k\equiv a^{\,k\bmod\varphi(m)}\pmod m\]
	\[k\ge\varphi(m)\Rightarrow
		a^k\equiv a^{\,k\bmod\varphi(m)+\varphi(m)}\pmod m\qquad(m\ge2)\]

\subsubsection{其他常用公式}
	$$\sum_{i = 1} ^ n \left[\gcd(i,n)=1\right]i
		=\frac n2\left(\varphi(n)+[n=1]\right)$$
	$$\sum_{i = 1} ^ n \sum_{j = 1} ^ i \left[\gcd(i,j)=d\right]
		=S_\varphi \left( \left\lfloor \frac n d \right\rfloor \right)$$
	$$\sum_{i = 1} ^ n \sum_{j = 1} ^ m \left[\gcd(i,j)=d\right]
		=\sum_{k=1}^{\left\lfloor\min(n,m)/d\right\rfloor}
		\mu(k)\left\lfloor\frac n{dk}\right\rfloor\left\lfloor\frac m{dk}\right\rfloor$$
	$$ \sum_{i = 1} ^ n f(i) \sum_{j = 1} ^ {\left\lfloor \frac n i \right\rfloor} g(j) = \sum_{i = 1} ^ n g(i) \sum_{j = 1} ^ {\left\lfloor \frac n i \right\rfloor} f(j) $$
\subsubsection{单位根反演}
\[
\omega_k=e^{2\pi\mathrm{i}/k},\qquad
\frac1k\sum_{j=0}^{k-1}\omega_k^{jr}=[k\mid r].
\]
\[
\begin{aligned}
A_t
&=\sum_{\substack{0\le i\le n\\i\equiv t\pmod k}}\binom ni
=\sum_{i=0}^n\binom ni[k\mid(i-t)]\\
&=\frac1k\sum_{j=0}^{k-1}\omega_k^{-jt}
	\sum_{i=0}^n\binom ni(\omega_k^j)^i\\
&=\boxed{\frac1k\sum_{j=0}^{k-1}
	\omega_k^{-jt}(1+\omega_k^j)^n}
\qquad(0\le t<k).
\end{aligned}
\]


\subsubsection{Arithmetic Function}
下文设 $m,n,u,v$ 为正整数，$f,g$ 为算术函数，
$n=\prod_{i=1}^{\omega(n)}p_i^{a_i}$（$p_i$ 为不同素数）；
$\tau(n)=\sum_{d\mid n}1$，$\omega(n)$ 为不同素因子数；
$*$ 表示 Dirichlet 卷积，$\cdot$ 表示逐点乘积。
\[p>1\text{ 时},\qquad p\text{ 为素数}\Longleftrightarrow
    (p-1)!\equiv-1\pmod p\]
\[a>1\Rightarrow
    \gcd(a^m-1,a^n-1)=a^{\gcd(m,n)}-1\]
\[\mu(n)^2=\sum_{d^2\mid n}\mu(d)\]
\[a>b>0,\ \gcd(a,b)=1\Rightarrow
    \gcd(a^m-b^m,a^n-b^n)=a^{\gcd(m,n)}-b^{\gcd(m,n)}\]
\[
P(m)=\prod_{\substack{1\le k\le m\\\gcd(k,m)=1}}k\qquad(m\ge2).
\]
\[
P(m)\equiv
\begin{cases}
-1,&m=4,p^q,2p^q,\\
1,&\text{其他情形}
\end{cases}
\pmod m,\qquad p\text{ 为奇素数},\ q\ge1.
\]
\[
\begin{aligned}
\sigma_k(n)&=\sum_{d\mid n}d^k
    =\prod_{i=1}^{\omega(n)}\sum_{j=0}^{a_i}p_i^{jk}&& (k\ge0),\\
\sigma_k(n)&=\prod_{i=1}^{\omega(n)}
    \frac{p_i^{(a_i+1)k}-1}{p_i^k-1}&& (k\ge1).
\end{aligned}
\]
\[J_k(n)=n^k\prod_{p\mid n}\left(1-\frac1{p^k}\right)\qquad(k\ge1)\]
$J_k(n)$ 等于满足 $1\le x_i\le n$ 且
$\gcd(x_1,\ldots,x_k,n)=1$ 的 $k$ 元组 $(x_1,\ldots,x_k)$ 数量。
\[\sum_{\delta\mid n}J_k(\delta)=n^k\]
\[\sum_{i=1}^n\sum_{j=1}^n[\gcd(i,j)=1]ij
    =\sum_{i=1}^n i^2\varphi(i)\]
\[\sum_{\delta\mid n}\delta^sJ_r(\delta)J_s\left(\frac n\delta\right)
    =J_{r+s}(n)\qquad(r,s\ge1)\]
\begin{align*}
\sum_{\delta\mid n}\varphi(\delta)\tau\left(\frac n\delta\right)=\sigma_1(n),&\quad
\sum_{\delta\mid n}|\mu(\delta)|=2^{\omega(n)},\\
\sum_{\delta\mid n}2^{\omega(\delta)}=\tau(n^2),&\quad
\sum_{\delta\mid n}\tau(\delta^2)=\tau(n)^2,\\
\sum_{\delta\mid n}\tau\left(\frac n\delta\right)2^{\omega(\delta)}=\tau(n)^2,&\quad
\sum_{\delta\mid n}\frac{\mu(\delta)}\delta=\frac{\varphi(n)}n,\\
\sum_{\delta\mid n}\frac{\mu(\delta)}{\varphi(\delta)}
    =\prod_{p\mid n}\frac{p-2}{p-1},&\quad
\sum_{\delta\mid n}\frac{\mu(\delta)^2}{\varphi(\delta)}
    =\frac n{\varphi(n)}.
\end{align*}
\[a\ge2\Rightarrow n\mid\varphi(a^n-1)\]
\[\sum_{\substack{1\le k\le n\\\gcd(k,n)=1}}f(\gcd(k-1,n))
    =\varphi(n)\sum_{d\mid n}\frac{(\mu*f)(d)}{\varphi(d)}\]
\[\varphi(\operatorname{lcm}(m,n))\varphi(\gcd(m,n))
    =\varphi(m)\varphi(n)\]
\[\sum_{\delta\mid n}\tau(\delta)^3
    =\left(\sum_{\delta\mid n}\tau(\delta)\right)^2\]
\[\tau(uv)=\sum_{\delta\mid\gcd(u,v)}\mu(\delta)
    \tau\left(\frac u\delta\right)\tau\left(\frac v\delta\right)\]
\[\sigma_k(u)\sigma_k(v)=\sum_{\delta\mid\gcd(u,v)}\delta^k
    \sigma_k\left(\frac{uv}{\delta^2}\right)\qquad(k\ge0)\]
\[\mu(n)=\sum_{k=1}^n[\gcd(k,n)=1]\cos\frac{2\pi k}n\]
\[\varphi(n)=\sum_{k=1}^n[\gcd(k,n)=1]
    =\sum_{k=1}^n\gcd(k,n)\cos\frac{2\pi k}n\]
\[
\left\{
\begin{aligned}
S(n)&=\sum_{k=1}^n(f*g)(k),\\
\sum_{k=1}^nS\left(\left\lfloor\frac nk\right\rfloor\right)
    &=\sum_{i=1}^nf(i)\sum_{j=1}^{\lfloor n/i\rfloor}(g*1)(j).
\end{aligned}
\right.
\]
\[
\left\{
\begin{aligned}
S(n)&=\sum_{k=1}^n(f\cdot g)(k),\quad g\text{ 完全积性},\\
\sum_{k=1}^nS\left(\left\lfloor\frac nk\right\rfloor\right)g(k)
    &=\sum_{k=1}^n(f*1)(k)g(k).
\end{aligned}
\right.
\]
\subsubsection{Binomial Coefficients}
组合恒等式中的整型变量均为非负整数，越界组合数取 $0$。
\begin{multicols*}{2}
\begin{scriptsize}
{\setlength{\tabcolsep}{0.75pt}
\begin{tabular}{|P{0.2cm}|P{0.3cm}|P{0.3cm}|P{0.3cm}|P{0.3cm}|P{0.3cm}|P{0.3cm}|P{0.3cm}|P{0.3cm}|P{0.3cm}|P{0.3cm}|P{0.3cm}|}
	\hline
	C & 0 & 1 & 2 & 3 & 4 & 5 & 6 & 7 & 8 & 9 & 10\\
	\hline
	0 & \cellcolor{black!20}1 &  &  &  &  &  &  &  &  &  & \\
	\hline
	1 & \cellcolor{black!20}1 & \cellcolor{black!20}1 &  &  &  &  &  &  &  &  & \\
	\hline
	2 & \cellcolor{black!20}1 & \cellcolor{black!0}2 & \cellcolor{black!20}1 &  &  &  &  &  &  &  & \\
	\hline
	3 & \cellcolor{black!20}1 & \cellcolor{black!20}3 & \cellcolor{black!20}3 & \cellcolor{black!20}1 &  &  &  &  &  &  & \\
	\hline
	4 & \cellcolor{black!20}1 & \cellcolor{black!0}4 & \cellcolor{black!0}6 & \cellcolor{black!0}4 & \cellcolor{black!20}1 &  &  &  &  &  & \\
	\hline
	5 & \cellcolor{black!20}1 & \cellcolor{black!20}5 & \cellcolor{black!0}10 & \cellcolor{black!0}10 & \cellcolor{black!20}5 & \cellcolor{black!20}1 &  &  &  &  & \\
	\hline
	6 & \cellcolor{black!20}1 & \cellcolor{black!0}6 & \cellcolor{black!20}15 & \cellcolor{black!0}20 & \cellcolor{black!20}15 & \cellcolor{black!0}6 & \cellcolor{black!20}1 &  &  &  & \\
	\hline
	7 & \cellcolor{black!20}1 & \cellcolor{black!20}7 & \cellcolor{black!20}21 & \cellcolor{black!20}35 & \cellcolor{black!20}35 & \cellcolor{black!20}21 & \cellcolor{black!20}7 & \cellcolor{black!20}1 &  &  & \\
	\hline
	8 & \cellcolor{black!20}1 & \cellcolor{black!0}8 & \cellcolor{black!0}28 & \cellcolor{black!0}56 & \cellcolor{black!0}70 & \cellcolor{black!0}56 & \cellcolor{black!0}28 & \cellcolor{black!0}8 & \cellcolor{black!20}1 &  & \\
	\hline
	9 & \cellcolor{black!20}1 & \cellcolor{black!20}9 & \cellcolor{black!0}36 & \cellcolor{black!0}84 & \cellcolor{black!0}126 & \cellcolor{black!0}126 & \cellcolor{black!0}84 & \cellcolor{black!0}36 & \cellcolor{black!20}9 & \cellcolor{black!20}1 & \\
	\hline
	10 & \cellcolor{black!20}1 & \cellcolor{black!0}10 & \cellcolor{black!20}45 & \cellcolor{black!0}120 & \cellcolor{black!0}210 & \cellcolor{black!0}252 & \cellcolor{black!0}210 & \cellcolor{black!0}120 & \cellcolor{black!20}45 & \cellcolor{black!0}10 & \cellcolor{black!20}1\\
	\hline
\end{tabular}}
\end{scriptsize}
\[ \binom{n}{k} \equiv [(n\mathbin{\&}k)=k] \pmod 2 \]
\[ \sum_{k=0}^n\binom{k}{m} = \binom{n+1}{m+1} \]
\[ \sqrt{1+z} = 1 + \sum_{k=1}^{\infty}\frac{(-1)^{k-1}}{k\times2^{2k-1}}\binom{2k-2}{k-1}z^k \]
\[ \begin{gathered}
\sum_{k=0}^{r}\binom{r-k}{m}\binom{s+k}{n}\\
=\binom{r+s+1}{m+n+1}\qquad(n\ge s)
\end{gathered} \]
{\scriptsize
\[ C_{n,m}:=\binom{n+m}{m}-\binom{n+m}{m-1}\qquad(n\ge m) \]
}
\end{multicols*}
\[ \binom{n}{k}=(-1)^k\binom{k-n-1}{k}
\quad\text{（右侧按广义组合数定义）},\qquad
\sum_{k=0}^{n}\binom{r+k}{k}=\binom{r+n+1}{n} \]
\[ \binom{n_1+\cdots+n_p}{m}
=\sum_{k_1+\cdots+k_p=m}\binom{n_1}{k_1}\cdots\binom{n_p}{k_p}\]
\subsubsection{Fibonacci Numbers, Lucas Numbers}
约定 $f_0=0,f_1=1$；整型变量非负，且各式所涉下标和运算均有定义。
\begin{multicols*}{2}
\[ F(z) = \frac{z}{1-z-z^2} \]
\[ \hat{\phi} = \frac{1-\sqrt{5}}{2} \]
\[ \sum_{k=1}^nf_k = f_{n+2}-1,\ \ \sum_{k=1}^nf^2_k = f_nf_{n+1} \]
\[ \sum_{k=0}^nf_kf_{n-k} = \frac{1}{5}(n-1)f_n+\frac{2}{5}nf_{n-1} \]
\[ \frac{f_{2n}}{f_n} = f_{n-1} + f_{n+1}\]
\[\begin{gathered}
	f_1+2f_2+3f_3+\cdots+nf_n \\
	=nf_{n+2}-f_{n+3}+2
\end{gathered}\]
\[ \gcd(f_m,f_n)=f_{\gcd(m,n)}\]
\[ f^2_n + (-1)^n = f_{n+1}f_{n-1} \]
\[ f_{n+k} = f_nf_{k+1} + f_{n-1}f_k \]
\[ f_{2n+1} = f^2_n+f^2_{n+1} \]
\[ (-1)^kf_{n-k} = f_{n}f_{k-1} - f_{n-1}f_{k} \]
\begin{small}
\inputminted{python3}{sections/NumberTheory/assets/tbr/fib.py}
\end{small}
\end{multicols*}
\[ \text{Modulo }f_n, f_{mn+r} \equiv \left\{
\begin{aligned}
&f_r,& m \bmod 4 = 0; \\
&(-1)^{r+1}f_{n-r},& m \bmod 4 = 1; \\
&(-1)^nf_r,& m \bmod 4 = 2; \\
&(-1)^{r+1+n}f_{n-r},& m \bmod 4 = 3.
\end{aligned}
\right.
\]
\[ L_0=2,L_1=1,L_n=L_{n-1}+L_{n-2}=(\frac{1+\sqrt{5}}{2})^n+(\frac{1-\sqrt{5}}{2})^n \]
\[ L(x)=\frac{2-x}{1-x-x^2} \]
若 $L_n$ 是素数，则 $n=0$、$n$ 是素数或 $n$ 是 $2$ 的幂；
其中 $L_0,L_4,L_8,L_{16}$ 是素数。
$$\phi = \frac {1 + \sqrt 5} 2, \; \hat\phi = \frac {1 - \sqrt 5} 2$$
$$F_n = \frac {\phi^n - {\hat\phi} ^ n} {\sqrt 5}, \; L_n = \phi ^ n + {\hat\phi} ^ n$$
$$\frac {L_n + F_n \sqrt 5} 2 = \left( \frac {1 + \sqrt 5} 2 \right) ^ n$$

\subsubsection{Sum of Powers}
\[\sum_{i=1}^ni^2=\frac{n(n+1)(2n+1)}{6},\ \ \sum_{i=1}^ni^3=\left(\frac{n(n+1)}{2}\right)^2\]
\[\sum_{i=1}^ni^4=\frac{n(n+1)(2n+1)(3n^2+3n-1)}{30}\]
\[\sum_{i=1}^ni^5=\frac{n^2(n+1)^2(2n^2+2n-1)}{12}\]

\subsubsection{Catalan Numbers 1, 1, 2, 5, 14, 42, 132, 429, 1430\dots}
\[c_{0}=1,c_n=\sum_{i=0}^{n-1}c_ic_{n-1-i}=c_{n-1}\frac{4n-2}{n+1}=\frac{\binom{2n}{n}}{n+1}=\binom{2n}{n}-\binom{2n}{n-1}\]
\[ c(x)=\frac{1-\sqrt{1-4x}}{2x}\]
Usage: $n$对括号序列; 含$n$个内部节点（$n+1$个叶子）的满二叉树;
$n\times n$的方格左下到右上不过对角线方案数; 凸$n+2$边形三角形分割数;
$n$个数的出栈方案数; 圆上$2n$个顶点两两连接且线段互不相交的方案数.
\paragraph*{类卡特兰数}
从 $(1,1)$ 出发走到 $(n,m)$,只能向右或者向上走,不能越过 $y=x$ 这条线(即保证 $x\geq y$ ). 当 $n\geq m$ 时,合法方案数是 $\binom{n+m-2}{m-1}-\binom{n+m-2}{m-2}$ (例如 $(3,2)$ 时为 $3-1=2$); 当 $n<m$ 时,合法方案数为 $0$.
\subsubsection{Motzkin Numbers 1, 1, 2, 4, 9, 21, 51, 127, 323, 835\dots}
圆上$n$点间选取若干条两两不相交弦的方案数。等价地，选
$k_1,\ldots,k_n\in\{-1,0,1\}$，使
$\sum_{i=1}^{a}k_i\ge0\ (1\le a\le n)$ 且 $\sum_{i=1}^{n}k_i=0$。
取 $M_0=M_1=1$；以下 $n\ge1$。
\[M_{n+1}=M_n+\sum_{i=0}^{n-1}M_iM_{n-1-i}=\frac{(2n+3)M_n+3nM_{n-1}}{n+3}\]
\[M_n=\sum_{k=0}^{\lfloor n/2\rfloor}\binom{n}{2k}c_k\]
\[M(x)=\frac{1-x-\sqrt{1-2x-3x^2}}{2x^2}\]
\subsubsection{Derangement 错排数 1, 0, 1, 2, 9, 44, 265, 1854, 14833\dots}
\[D_0=1,\qquad D_n=n!\sum_{k=0}^{n}\frac{(-1)^k}{k!}\]
\[D_n=(n-1)(D_{n-1}+D_{n-2})\qquad(n\ge2)\]
\subsubsection{Bell Numbers 1, 1, 2, 5, 15, 52, 203, 877, 4140 \dots}
$n$个元素集合划分的方案数.
\[ B_n=\sum_{k=0}^{n}\stirlingtwo{n}{k},\ \ B_{n+1} = \sum_{k=0}^n\binom{n}{k}B_k \]
\[ B_{p^m+n} \equiv mB_n+B_{n+1} \pmod{p}
\qquad(p\text{ 为素数},\ m,n\ge0) \]
\[B(x)=\sum_{n=0}^{\infty}\frac{B_n}{n!}x^n=\mathrm{e}^{\mathrm{e}^x-1}\]
\subsubsection{Stirling Numbers}
约定越界的 Stirling 数为 $0$。
\paragraph*{第一类} $n$个元素集合分作$k$个\textbf{非空轮换}方案数.
\begin{multicols*}{2}
$$\begin{aligned}
\stirlingone{n+1}{k} =& n\stirlingone{n}{k} + \stirlingone{n}{k-1} \\
s(n,k) &=(-1)^{n-k}\stirlingone{n}{k}
\end{aligned}$$
\begin{tabular}{|P{0.39cm}|P{0.39cm}|P{0.45cm}|P{0.57cm}|P{0.57cm}|P{0.45cm}|P{0.39cm}|P{0.39cm}|}
	\hline
	n\textbackslash k & 0 & 1 & 2 & 3 & 4 & 5 & 6 \\
	\hline
	0 & 1 &  &  &  &  &  &  \\
	\hline
	1 & 0 & 1 &  &  &  &  &  \\
	\hline
	2 & 0 & 1 & 1 &  &  &  &  \\
	\hline
	3 & 0 & 2 & 3 & 1 &  &  &  \\
	\hline
	4 & 0 & 6 & 11 & 6 & 1 &  &  \\
	\hline
	5 & 0 & 24 & 50 & 35 & 10 & 1 &  \\
	\hline
	6 & 0 & 120 & 274 & 225 & 85 & 15 & 1 \\
	\hline
	7 & 0 & 720 & 1764 & 1624 & 735 & 175 & 21 \\
	\hline
\end{tabular}

$$\begin{aligned}
	\stirlingone{n+1}{2} &= n!H_n \text{ (see \ref{section:Harmonic})} \\
	x^{\underline{n}} &= \sum_{k=0}^{n} \stirlingone{n}{k}(-1)^{n-k}x^k \\
	x^{\overline{n}} &= \sum_{k=0}^{n} \stirlingone{n}{k}x^k
\end{aligned}$$

For fixed $k$, EGF:

\[\begin{aligned}
	\sum_{n = 0} ^ \infty \stirlingone{n}{k} \frac {x ^ n} {n!}
	&= \frac {x ^ k} {k!} \left( \frac {-\ln (1 - x)} x \right) ^ k \\
	&\quad (k=1:\,-\ln(1-x))
\end{aligned}\]
\end{multicols*}

\paragraph*{第二类} 把$n$个元素集合分作$k$个\textbf{非空子集}方案数.

\begin{multicols*}{2}
$$\begin{aligned}
\stirlingtwo{n+1}{k} &= k\stirlingtwo{n}{k} + \stirlingtwo{n}{k-1} \\
m!\stirlingtwo{n}{m} &= \sum_{k=0}^{m}\binom{m}{k}k^n(-1)^{m-k}
\end{aligned}$$
\begin{tabular}{|P{0.39cm}|P{0.39cm}|P{0.39cm}|P{0.39cm}|P{0.45cm}|P{0.45cm}|P{0.45cm}|P{0.39cm}|}
	\hline
	n\textbackslash k & 0 & 1 & 2 & 3 & 4 & 5 & 6 \\
	\hline
	0 & 1 &  &  &  &  &  &  \\
	\hline
	1 & 0 & 1 &  &  &  &  &  \\
	\hline
	2 & 0 & 1 & 1 &  &  &  &  \\
	\hline
	3 & 0 & 1 & 3 & 1 &  &  &  \\
	\hline
	4 & 0 & 1 & 7 & 6 & 1 &  &  \\
	\hline
	5 & 0 & 1 & 15 & 25 & 10 & 1 &  \\
	\hline
	6 & 0 & 1 & 31 & 90 & 65 & 15 & 1 \\
	\hline
	7 & 0 & 1 & 63 & 301 & 350 & 140 & 21 \\
	\hline
\end{tabular}
\columnbreak
$$\begin{aligned}
	x^n &= \sum_{k=0}^{n} \stirlingtwo{n}{k}x^{\underline{k}} \\
	&= \sum_{k=0}^{n} \stirlingtwo{n}{k}(-1)^{n-k}x^{\overline{k}}
\end{aligned}$$
For fixed $k$, EGF and OGF:
$$\begin{aligned}
\sum_{n = 0} ^ \infty \stirlingtwo{n}{k} \frac {x ^ n} {n!} &= \frac {x ^ k} {k!} \left( \frac {e ^ x - 1} x \right) ^ k \\
\sum_{n = 0} ^ \infty \stirlingtwo{n}{k} x ^ n &= x ^ k \prod_{i = 1} ^ k (1 - i x) ^ {-1}
\end{aligned}$$
\end{multicols*}
\subsubsection{Eulerian Numbers}
以下 $n\ge1$，显式公式中 $0\le m<n$；约定
$\bangle{n}{k}=0$ 当 $k\notin[0,n-1]$。
\begin{multicols*}{2}
	{\setlength{\tabcolsep}{0.75pt}
	\begin{tabular}{|P{0.5cm}|P{0.3cm}|P{0.5cm}|P{0.57cm}|P{0.57cm}|P{0.57cm}|P{0.5cm}|P{0.3cm}|}
		\hline
		n\textbackslash k & 0 & 1 & 2 & 3 & 4 & 5 & 6\\
		\hline
		1 & 1 &  &  &  &  &  & \\
		\hline
		2 & 1 & 1 &  &  &  &  & \\
		\hline
		3 & 1 & 4 & 1 &  &  &  & \\
		\hline
		4 & 1 & 11 & 11 & 1 &  &  & \\
		\hline
		5 & 1 & 26 & 66 & 26 & 1 &  & \\
		\hline
		6 & 1 & 57 & 302 & 302 & 57 & 1 & \\
		\hline
		7 & 1 & 120 & 1191 & 2416 & 1191 & 120 & 1\\
		\hline
	\end{tabular}}
	\columnbreak
		\[ \bangle{n}{k} = (k+1)\bangle{n-1}{k} + (n-k)\bangle{n-1}{k-1} \]
		\[ x^n = \sum_{k=0}^{n-1} \bangle{n}{k}\binom{x+k}{n} \]
		\[ \bangle{n}{m} = \sum_{k=0}^m\binom{n+1}{k}(m+1-k)^n(-1)^k \]
\end{multicols*}
\begin{multicols*}{2}
\subsubsection{Harmonic Numbers, 1, 3/2, 11/6, 25/12, 137/60\dots}
\label{section:Harmonic}
以下 $n\ge1$，且第三式中 $0\le m\le n$。
\[ H_n = \sum_{k=1}^n \frac{1}{k}, \sum_{k=1}^nH_k = (n+1)H_n-n \]
\[ \sum_{k=1}^nkH_k = \frac{n(n+1)}{2}H_n - \frac{n(n-1)}{4} \]
\[ \sum_{k=1}^n\binom{k}{m}H_k = \binom{n+1}{m+1}(H_{n+1} - \frac{1}{m+1}) \]

\subsubsection{卡迈克尔函数}
卡迈克尔函数表示模 $m$ 简化剩余系的指数，即
\[\lambda(m)=\max_{(a,m)=1}\delta_m(a).\]

当 $m\ge2$ 时，$m$ 存在原根当且仅当 $\lambda(m)=\varphi(m)$。

若 $m=\prod_{i=1}^t p_i^{\alpha_i}$，则
\[\lambda(m)=\operatorname{lcm}_{1\le i\le t}\lambda(p_i^{\alpha_i}).\]
奇质数 $p$ 满足 $\lambda(p^\alpha)=(p-1)p^{\alpha-1}$；此外
\[\begin{gathered}
\lambda(1)=\lambda(2)=1,\qquad \lambda(4)=2,\\
\lambda(2^\alpha)=2^{\alpha-2}\quad(\alpha\ge3).
\end{gathered}\]

\end{multicols*}

\subsubsection{求拆分数}
\begin{multicols*}{2}
\inputminted{python3}{sections/NumberTheory/assets/yzh/partition.py}
\[\begin{aligned}
\Phi(x)&=\prod_{j=1}^{\infty}(1-x^j)\\
&=\sum_{t=-\infty}^{\infty}(-1)^t x^{t(3t-1)/2}.
\end{aligned}\]

\newcolumn
	记 $p(n)$ 为 $n$ 的拆分数，$f(n,k)$ 为每种正整数至多使用 $k$ 次的拆分数，并令相应生成函数为 $P(x),F_k(x)$。则
	\[\begin{aligned}
	P(x)&=\frac1{\Phi(x)},\\
	F_k(x)&=\frac{\Phi(x^{k+1})}{\Phi(x)}.
	\end{aligned}\]
	约定 $p(0)=1,p(t)=0\ (t<0)$。由五边形数定理可在 $O(n\sqrt n)$ 内求出前 $n$ 项：\\
	$p(n) = p(n-1)+p(n-2)-p(n-5)-p(n-7)+\cdots$, \\
	$f(n,k)=p(n)-p(n-k-1)-p(n-2k-2)+p(n-5k-5)+p(n-7k-7)-\cdots$。
\end{multicols*}

\subsubsection{Bernoulli Numbers 1, -1/2, 1/6, 0, -1/30, 0, 1/42 \dots}

$$\begin{aligned}B(x)=\sum_{i\ge 0}\frac{B_i x^i}{i!}=\frac x{e^x-1}\end{aligned}$$

\[B_0=1,\qquad B_n=-\frac1{n+1}\sum_{i=0}^{n-1}\binom{n+1}{i}B_i\quad(n\ge1).\]
\[\sum_{i=0}^n\binom{n+1}{i}B_i=0\quad(n\ge1).\]

$$ \begin{aligned}S_n(m)=\sum_{i=0}^{m-1}i^n=\sum_{i=0}^n\binom{n}{i}B_{n-i}\frac{m^{i+1}}{i+1}\end{aligned} $$

$ B_0 = 1,\; B_1 = -\frac 1 2,\; B_2=\frac16,\; B_4 = -\frac 1 {30},\; B_6 = \frac 1 {42},\; B_8 = -\frac 1{30},\; \dots$.

(除了$B_1 = -\frac 1 2$以外, 伯努利数的奇数项都是$0$.)

自然数幂次和关于次数的EGF:
\[ F_n(x)=\sum_{k=0}^\infty \frac{\sum_{i=0}^n i^k}{k!}x^k
=\sum_{i=0}^n e^{ix}=\frac{e^{(n+1)x}-1}{e^x-1}. \]

\subsubsection{kMAX-MIN反演}
设 $k\mathrm{MAX}(S)$ 为有限集合 $S$ 中第 $k$ 大元素，其中
$1\le k\le |S|$；约定组合数越界时为 $0$。
\[k \mathrm{MAX}(S)=\sum_{T\subset S, T\neq \emptyset}(-1)^{|T|-k}C_{|T|-1}^{k-1}\mathrm{MIN}(T)\]
代入$k=1$即为MAX-MIN反演:
\[\mathrm{MAX}(S)=\sum_{T\subset S, T\neq \emptyset}(-1)^{|T|-1}\mathrm{MIN}(T)\]
\subsubsection{伍德伯里矩阵恒等式}
设 $A\in\mathbb F^{n\times n},U\in\mathbb F^{n\times r},
C\in\mathbb F^{r\times r},V\in\mathbb F^{r\times n}$，且下式中的逆均存在，则
$$(A+UCV)^{-1}=A^{-1}-A^{-1}U(C^{-1}+VA^{-1}U)^{-1}VA^{-1}$$
秩一修改时取 $r=1,C=[1]$。修改第 $i$ 行可令 $U=e_i$，修改第 $j$ 列可令 $V=e_j^\mathsf T$；一次维护复杂度为 $O(n^2)$。
\subsubsection{Sum of Squares}
$r_k(n)=\#\{(x_1,\ldots,x_k)\in\mathbb Z^k:\sum_i x_i^2=n\}$。设
\[n=2^{a_0}\prod_{i=1}^r p_i^{a_i}\prod_{j=1}^s q_j^{b_j},\qquad
p_i\equiv3\pmod4,\ q_j\equiv1\pmod4.\]
那么
\[r_2(n)=\left\{\begin{aligned}
&0, & \text{某个 }a_i\text{ 为奇数},\\
&4\prod_{j=1}^s(b_j+1), & \text{所有 }a_i\text{ 均为偶数}.
\end{aligned}\right.\]
$r_3(n)>0$ 当且仅当 $n$ 不能写成 $4^a(8b+7)$，其中 $a,b\ge0$ 为整数。
\subsubsection{枚举勾股数 Pythagorean Triple}
所有正整数勾股数组（允许交换 $x,y$）可写成
\[x=d(m^2-n^2),\quad y=2dmn,\quad z=d(m^2+n^2),\]
其中 $d\ge1,m>n\ge1$。它是本原勾股数组当且仅当
$d=1,(m,n)=1$ 且 $m-n$ 为奇数；枚举全部勾股数组时还需枚举倍数 $d$。
\subsubsection{四面体体积 Tetrahedron Volume}
设 $U,V,W$ 构成四面体的一个面，$u,v,w$ 分别为它们的对边（$u$ 与 $U$ 相对，以此类推），并同步轮换两组字母，则体积为
\[ \mathcal V = \frac{\sqrt{ 4u^2v^2w^2 - \sum_{\mathrm{cyc}}{u^2(v^2+w^2-U^2)^2} + \prod_{\mathrm{cyc}}{(v^2+w^2-U^2)} }}{12} \]
\subsubsection{杨氏矩阵 与 钩子公式}
形状 $\lambda$ 的标准杨表将 $1,\ldots,n$ 各填一次，向右、向上严格递增，其中 $n=|\lambda|$。令 $F_n$ 为大小为 $n$ 的所有标准杨表数之和，则
\[F_0=F_1=1,\quad F_n=F_{n-1}+(n-1)F_{n-2}\ (n\ge2),\qquad
\sum_{n\ge0}F_n\frac{x^n}{n!}=e^{x+x^2/2}.\]
钩长公式给出形状 $\lambda$ 的标准杨表数：
\[d_\lambda=\frac{n!}{\prod_{(i,j)\in\lambda}h_\lambda(i,j)}.\]
$h_\lambda(i,j)$ 为该格右边、上边的格子数之和再加 $1$。

\subsubsection{常见博弈游戏}
\begin{multicols}{2}
%\paragraph{减法游戏}$n$根火柴每次可以拿$a_i$根, 最后一个拿的人获胜. 
%\paragraph{Ferguson游戏}两个盒子, 一个有$m$个糖另一个有$n$个糖, 每次走步将一个盒子清空而把另一个盒子的一些糖拿到被清空的盒子中, 使得两个盒子各至少有一颗糖. 到终态(1,1)的人获胜. 
%\paragraph{动态减法游戏}初始有$n$根火柴, 第一个人可以任意取, 至少$1$根至多$n-1$根, 以后每个人至少拿一根, 但拿的火柴数不能超过上一次拿的火柴数.
\paragraph{Nim-K游戏}正常规则下有 $n$ 堆石子，每次从 $1$ 至 $k$ 个非空堆中各取正数颗，取走最后一颗者胜。将各堆大小二进制分解，先手必败当且仅当每一位上 $1$ 的个数均为 $k+1$ 的倍数。
\paragraph{Anti-Nim游戏}每次只操作一堆，取走最后一颗者输。若所有非空堆均为 $1$，先手胜当且仅当堆数为偶数；否则先手胜当且仅当各堆大小的普通 Nim 异或和非零。
\paragraph{斐波那契博弈}一堆有 $n$ 件物品，首回合取 $1$ 至 $n-1$ 件，以后每次取 $1$ 至上次取数的两倍，取走最后一件者胜。先手胜当且仅当 $n$ 不是斐波那契数。
\paragraph{威佐夫博弈}两堆石子，每次从一堆取正数颗，或从两堆各取相同的正数颗，取完者胜。令 $C=|A-B|$，则后手胜当且仅当 $\left\lfloor C\frac{\sqrt5+1}{2}\right\rfloor=\min(A,B)$。
\paragraph{约瑟夫环}$n$ 个人编号 $0,\dots,n-1$，令 $f_{i,m}$ 为 $i$ 个人时报到 $m$ 被淘汰规则下的最终胜者，则 $f_{1,m}=0$，$f_{i,m}=(f_{i-1,m}+m)\bmod i\ (i\ge2)$。
\paragraph{阶梯Nim}台阶编号为 $0,1,\dots,n$，$0$ 为终点；每次将某个 $i\ge1$ 上的正数颗石子移到 $i-1$，最后移动者胜。SG 值为奇数编号台阶石子数的异或和。树上版本以根为终点、每次移至父亲，结论为奇数深度节点石子数的异或和。
\paragraph{图上博弈}给定无向图和起点 $v$，双方轮流走到相邻且未访问的点，无法移动者输。设 $\nu(G)$ 为最大匹配大小，则先手败当且仅当 $\nu(G-v)=\nu(G)$，即存在不覆盖 $v$ 的最大匹配。
\end{multicols}
\paragraph{极小极大定理与零和博弈均衡}
设行玩家选 $i$、列玩家选 $j$ 时，列玩家收益为 $w_{i,j}$。双方的对偶线性规划为
\[\begin{array}{c@{\qquad}c}
\begin{aligned}
\max_{\mathbf p,v}\quad &v\\[-2pt]
\mathrm{s.t.}\quad &\sum_jw_{i,j}p_j\ge v\ (\forall i),\\[-2pt]
&p_j\ge0,\ \sum_jp_j=1
\end{aligned}
&
\begin{aligned}
\min_{\mathbf q,u}\quad &u\\[-2pt]
\mathrm{s.t.}\quad &\sum_iq_iw_{i,j}\le u\ (\forall j),\\[-2pt]
&q_i\ge0,\ \sum_iq_i=1.
\end{aligned}
\end{array}\]
极小极大定理保证 $u^*=v^*$；两边的最优策略共同构成纳什均衡。只解其中一边只能得到该玩家的最优策略，一般应使用线性规划而非贪心调整概率。


\subsubsection{概率相关}
\[\begin{aligned}
D(X)&=E[(X-E[X])^2]=E[X^2]-E[X]^2,\quad D(aX)=a^2D(X),\\
D(X+Y)&=D(X)+D(Y)+2\operatorname{Cov}(X,Y).
\end{aligned}\]
若 $X,Y$ 不相关，则协方差项为 $0$；若 $X$ 为非负整数值随机变量，则
\[E[X]=\sum_{i=1}^{\infty}P(X\ge i).\]
$m$ 个数的总体方差为 $s^2=m^{-1}\sum_{i=1}^m x_i^2-\overline x^2$。

\subsubsection{邻接矩阵行列式的意义}
对 $n$ 阶邻接矩阵 $A$，
\[\det A=\sum_{\sigma\in S_n}\operatorname{sgn}(\sigma)
\prod_{i=1}^n A_{i,\sigma(i)}.\]
非零项对应覆盖全部顶点的有向循环覆盖。若置换 $\sigma$ 含 $e(\sigma)$ 个偶长度循环，则 $\operatorname{sgn}(\sigma)=(-1)^{e(\sigma)}$。对无向对称矩阵，长度至少为 $3$ 的环有两个方向；二元环 $u\leftrightarrow v$ 贡献 $A_{uv}A_{vu}$，对称边权时即 $w_{uv}^2$。

\subsubsection{Others (某些近似数值公式在这里)}
\[ S_j = \sum_{k=1}^nx_k^j \]
\[ e_m = \sum_{1\leq j_1 < \cdots < j_m \leq n} x_{j_1}\cdots x_{j_m},\ \ h_m = \sum_{1\leq j_1 \leq \cdots \leq j_m \leq n} x_{j_1}\cdots x_{j_m} \]
约定 $e_0=h_0=1$。对 $m\ge1$，Newton 恒等式为
\[\begin{aligned}
e_m&=\frac1m\sum_{j=1}^m(-1)^{j+1}S_je_{m-j},\\
h_m&=\frac1m\sum_{j=1}^mS_jh_{m-j}.
\end{aligned}\]
\[ \sum_{k=0}^nkc^k = \frac{nc^{n+2}-(n+1)c^{n+1}+c}{(c-1)^2}\quad(c\ne1), \qquad
\sum_{k=0}^n k=\frac{n(n+1)}2\quad(c=1). \]
\[\begin{aligned}
H_n&=\sum_{i=1}^n\frac1i\\
&=\ln\left(n+\frac12\right)+\gamma+\frac{1}{24(n+1/2)^2}+O(n^{-4}).
\end{aligned}\]
$\gamma\approx0.5772156649015328606$ 为欧拉常数。
\[ n! = \sqrt{2\pi n}\left(\frac{n}{e}\right)^n
\left(1+\frac{1}{12n}+\frac{1}{288n^2}+O(n^{-3})\right) \]
\[ \begin{aligned}
 &\max{\{x_a-x_b, y_a-y_b, z_a-z_b\}} - \min{\{x_a-x_b, y_a-y_b, z_a-z_b\}} \\
=& \frac{1}{2}\sum_{\mathrm{cyc}}\left| (x_a-y_a)-(x_b-y_b) \right|
\end{aligned} \]
\[ (a+b)(b+c)(c+a) = \frac{(a+b+c)^3 - a^3 - b^3 - c^3}{3} \]
\[ a^3+b^3=(a+b)(a^2-ab+b^2),a^3-b^3=(a-b)(a^2+ab+b^2) \]
$n\ge1$ 为奇数时：
\[ a^n+b^n=(a+b)(a^{n-1}-a^{n-2}b+a^{n-3}b^2-\cdots-ab^{n-2}+b^{n-1})\]
$n$ 个一般位置的仿射超平面最多把 $k$ 维空间分成
$\sum_{i=0}^{k}\binom ni$ 个区域。
\end{footnotesize}
